
\section{Requirements}
In the context of software, requirements state boundaries or necessities of specific circumstances or events that occur, so that the software is still able to run without any inconvenience. Requirements in the field of product ownership generally means what exactly the software does and what not. The requirements get communicated mainly with the customers and the product owner, who is trained to figure out the exact needs for the software.
Below, the general requirements for the system are first defined in functional and non-functional terms. Then, the specific requirements for the core components are presented in detail.

\subsection{Requirements}

\renewcommand{\arraystretch}{1.3}
% Flattersatz für p{}-Spalten
\newcolumntype{P}[1]{>{\raggedright\arraybackslash}p{#1}}

\begin{longtable}{|P{1.5cm}|P{4cm}|P{8cm}|P{2cm}|}
\hline
\rowcolor{gray!20}
\multicolumn{4}{|l|}{\textbf{1. Overall System Requirements}} \\
\hline
\textbf{ID} & \textbf{Name} & \textbf{Description} & \textbf{Priority} \\
\hline
R1.1 & OSM-Integration & The system must be able to integrate maps from OpenStreetMap. & High \\
\hline
R1.2 & Automation & The system must enable the developed workflow (network \& traffic generation, simulation, analysis, output) to be automated and script-controlled. & High \\
\hline
R1.3 & Configurable Parameters & Users have to be able to set parameters like charging power, number of charging stations and agents, mobility patterns, charging behavior profiles and optionally power grid data. & Medium \\
\hline
R1.4 & Modularity & The system should have a modular structure so that components such as simulation, traffic generation and visualization can be processed independently. & Low \\
\hline
R1.5 & Extensibility & The system architecture should support future extensions (e.g. new optimization algorithms, data inputs). & Low \\
\hline
R1.6 & Performance & The system should handle large geographic areas and thousands of agents in $<$ 10  minutes. & High \\
\hline
R1.7 & Usability & The web UI should be intuitive, with real-time feedback and helpful defaults. & Medium \\
\hline
R1.8 & Reproducibility & The system must allow reproducibility via random seeds, logs, and versioning. & High \\
\hline
\rowcolor{gray!20}
\multicolumn{4}{|l|}{\textbf{2. Traffic and Charging Behavior}} \\
\hline
R2.1 & EVs with Charging Profiles & The system must simulate electric vehicles with energy-related behavior (SOC, charging-demand, consumption). & High \\
\hline
R2.2 & Charging Station Placement & It must be possible to define charging points in the network via script and independently of main traffic. & High \\
\hline
R2.3 & Bottleneck Detection & The system should detect bottlenecks at charging points. & Medium \\
\hline
R2.4 & Traffic Agent Generation & The system should generate realistic agents with daily routines (e.g. work-home-leisure), including rush hours. & High \\
\hline
R2.5 & Vehicle Configuration & It must support agent density, battery sizes, vehicle types. & Medium \\
\hline
\rowcolor{gray!20}
\multicolumn{4}{|l|}{\textbf{3. Location Analysis and Optimization}} \\
\hline
R3.1 & Clustering Analysis & Hotspot analysis for EV charging demand. & High \\
\hline
R3.2 & New Charging Point Proposals & Automatic suggestion of new charging locations. & High \\
\hline
R3.3 & Data Export & Export of results (GeoJSON, XML). & Medium \\
\hline
R3.4 & Multi-Criteria Optimization & Optimize locations based on traffic flow, accessibility, grid capacity. & High \\
\hline
R3.5 & Optimization Goals & Support goals like maximizing coverage or minimizing average distance. & Medium \\
\hline
R3.6 & Result Scoring & Provide ranked list of candidate sites with scores. & Medium \\
\hline
R3.7 & Logs & Maintain detailed logs for all optimization runs. & Medium \\
\hline
\rowcolor{gray!20}
\multicolumn{4}{|l|}{\textbf{4. Network Modeling and Geo-data Processing}} \\
\hline
R4.1 & Automatic conversion & Convert OpenStreetMap into simulation-ready format. & High \\
\hline
R4.2 & POIs & Automatically detect POIs from OSM data. & High \\
\hline
R4.3 & Network & Automatically generate simulation-ready network (.xml). & High \\
\hline
R4.4 & Grid Data Support & Optionally integrate power grid data into optimization. & Medium \\
\hline
\rowcolor{gray!20}
\multicolumn{4}{|l|}{\textbf{5. Population Behavior}} \\
\hline
R5.1 & Daily Patterns & The population must follow a reasonable daily pattern (home-work-leisure-home-patterns). & High \\
\hline
R5.2 & Temporal Variation & Include stochastic variation in departure/arrival times. & Medium \\
\hline
R5.3 & Population & Automatically generate simulation-ready population (\texttt{.pop}). & High \\
\hline
\rowcolor{gray!20}
\multicolumn{4}{|l|}{\textbf{6. Visualization and Validation}} \\
\hline
R6.1 & SUMO Integration & Calculated stations must be integrated in SUMO for validation. & High \\
\hline
R6.2 & Simulation Metrics & Collect charging behavior, grid load, waiting times. & High \\
\hline
R6.3 & Interactive Dashboard & Display station locations, KPIs, simulation results, export (CSV, GeoJSON). & Medium \\
\hline
\end{longtable}

\subsection{Success Criteria}

The system is successful if users can:
\begin{enumerate}
    \item Import arbitrary OSM data and configure parameters easily.
    \item Generate realistic traffic patterns with configurable parameters.
    \item Find charging station locations that perform well in simulation.
    \item Reproduce results consistently across multiple runs.
\end{enumerate}

The ultimate goal is providing a reliable research and planning tool that helps optimizing EV charging infrastructure placement through data-driven simulation and analysis.